% TOP_PROBLEM: {"problemId": "problem.vaught-s-conjecture", "canonicalTitle": "Vaught's conjecture", "releaseRank": 242, "primaryDomain": "logic_foundations_set_theory", "primaryDomainLabel": "Logic, foundations & set theory", "displayStatus": "open_with_solved_subcases", "releaseStatus": "open_with_solved_subcases"}
% !TeX program = lualatex
% Definition and sources only; this file is not an LLM solution attempt.
\documentclass[11pt]{article}
\usepackage[margin=25mm]{geometry}
\usepackage{fontspec}
\setmainfont{Latin Modern Roman}
\usepackage{amsmath,amssymb,mathtools}
\usepackage{unicode-math}
\setmathfont{Latin Modern Math}
\usepackage{xurl,hyperref}
\hypersetup{colorlinks=true,urlcolor=blue}
\setcounter{secnumdepth}{0}
\setlength{\parindent}{0pt}
\setlength{\parskip}{5pt}
\setlength{\emergencystretch}{3em}
\begin{document}
\section*{\texorpdfstring{Vaught's conjecture}{Vaught's conjecture}}\label{242}
\subsection{Definitions and mathematical statement}
For a complete first-order theory $T$ in a countable language, let $I(T,\aleph_0)$ be the number of isomorphism classes of its countable models. Completeness means that for each sentence, $T$ decides it or its negation. Vaught's assertion is
\[
 I(T,\aleph_0)\le\aleph_0\quad\text{or}\quad I(T,\aleph_0)=2^{\aleph_0}.
\]
Thus no intermediate uncountable number of countable models is permitted. Models are compared by isomorphism, not elementary equivalence, which all models of $T$ already share. The statement is about arbitrary complete countable theories; a result for one relational signature or a selected structural class is only a subcase.
\par\smallskip\textit{Formulation and concept references:} \href{https://www.sciencedirect.com/science/article/abs/pii/S0168007224000083}{[S1]}.

\subsection{Short English statement}
Must a complete theory in a countable language have either at most countably many countable models or as many as there are real numbers?
\subsection{Sources}
\begin{itemize}
\item[S1]Sharp Vaught's conjecture for some classes of partial orders.
\url{https://www.sciencedirect.com/science/article/abs/pii/S0168007224000083}.
\textit{Formulation source.}
\end{itemize}
\end{document}
